\section*{Alignment is not a Curve}

\begin{insight}
In most models, a railway alignment is described as a curve.

A mapping
\[
\gamma : s \mapsto (x(s), y(s))
\]
is considered sufficient.

\medskip

This view is convenient.

It is also incomplete.

\medskip

A curve describes position.

A railway alignment governs motion.

\medskip

For a train:
\begin{itemize}
\item curvature determines lateral acceleration,
\item curvature variation determines jerk,
\item and both directly affect safety, comfort, and wear.
\end{itemize}

\medskip

These quantities are not properties of position.

They are properties of change.

\medskip

Thus, the essential object is not the curve itself,
but the evolution of curvature.

\medskip

An alignment is not a static geometric shape.

It is a controlled process.

\medskip

\[
\text{alignment} \;\equiv\; \kappa(s)
\]

\medskip

The curve is only the result.

\medskip

To model alignment means to model how geometry changes —
not just where it is.
\end{insight}
